\section{Results}
Running the evaluation pipeline produced a comprehensive dataset with over 2,000 AI scoring judgments. Table \ref{tab:results} presents the aggregated evaluation results across all five experimental configurations, including Cohen's Kappa scores and percentage distributions for the four evaluation labels.

\begin{table}[H]
\centering
\caption{Aggregated RAG Evaluation Results across 416 Semantic Mutations and 5 Configurations.}
\label{tab:results}
\begin{tabular}{lccccc}
\toprule
\textbf{Model Configuration} & \textbf{Kappa ($k$)} & \textbf{Abstain} & \textbf{Faithful} & \textbf{Inconsistent} & \textbf{Hallucination} \\
\midrule
Llama 3.1 17B (With Prompt, Private Data) & 0.918 & 51.92\% & 45.19\% & 2.88\% & 0.00\% \\
Llama 3.1 17B (With Prompt) & 0.738 & 60.10\% & 30.53\% & 9.37\% & 0.00\% \\
Llama 3.1 17B (No Prompt) & 0.825 & 3.61\% & 51.44\% & 44.71\% & 0.24\% \\
Llama 3.1 8B (With Prompt) & 0.901 & 63.70\% & 30.29\% & 5.77\% & 0.24\% \\
Llama 3.1 8B (No Prompt) & 0.806 & 42.55\% & 27.64\% & 29.81\% & 0.00\% \\
\bottomrule
\end{tabular}
\end{table}

The Binomial tests on the Abstain rates yielded $p$-values approaching zero (e.g., $p = 5.76 \times 10^{-58}$ for the 17B With Prompt configuration, and $p = 5.61 \times 10^{-126}$ for the 8B No Prompt configuration). In every case, we strongly reject the null hypothesis ($H_0$), confirming that none of the tested configurations successfully maintained the $\ge 90\%$ defensive threshold.

To better understand the effect of each variable (model size, prompt engineering, and data source), we break down the results for each configuration below.

\subsection{Configuration 1: Llama 3.1 17B (With Prompt) -- Public Data}
\begin{table}[H]
\centering
\caption{Evaluation Results for Llama 3.1 17B (With Prompt). $k = 0.7375$.}
\label{tab:17b_prompt}
\begin{tabular}{lcc}
\toprule
\textbf{Classification} & \textbf{Count} & \textbf{Rate (\%)} \\
\midrule
Abstain & 250 & 60.10\% \\
Faithful & 127 & 30.53\% \\
Inconsistent & 39 & 9.37\% \\
Hallucination & 0 & 0.00\% \\
\bottomrule
\end{tabular}
\end{table}
This is considered the strongest baseline configuration in the study, using the 17B model combined with a strict defensive System Prompt, tested on the Public dataset (FastAPI documentation). The Abstain rate reached 60.10\%, showing that the System Prompt is fairly effective at steering model behavior when dealing with corrupted data. However, this number still falls short of the absolute safety threshold ($\ge 90\%$). Notably, 30.53\% of the time the model still passively accepted and followed the misinformation in the context (Faithful), suggesting that lexical instructions alone are not enough to fully eliminate the model's dependence on the provided context.

\subsection{Configuration 2: Llama 3.1 17B (No Prompt) -- Public Data}
\begin{table}[H]
\centering
\caption{Evaluation Results for Llama 3.1 17B (No Prompt). $k = 0.8246$.}
\label{tab:17b_noprompt}
\begin{tabular}{lcc}
\toprule
\textbf{Classification} & \textbf{Count} & \textbf{Rate (\%)} \\
\midrule
Faithful & 214 & 51.44\% \\
Inconsistent & 186 & 44.71\% \\
Abstain & 15 & 3.61\% \\
Hallucination & 1 & 0.24\% \\
\bottomrule
\end{tabular}
\end{table}
Removing the System Prompt led to a dramatic collapse of the defensive line. The 17B model almost never refused to answer (Abstain dropped to just 3.61\%). Instead, thanks to its strong reading comprehension, the model's behavior split into two opposite tendencies: either it fully trusted the corrupted text (Faithful shot up to 51.44\%), or it automatically drew on its deep background knowledge to correct the misinformation (Inconsistent reached 44.71\%). It is worth noting that the model correcting information on its own --- even if the corrections are accurate --- actually violates the groundedness principle, which is a core requirement for any RAG system.

\subsection{Configuration 3: Llama 3.1 8B (With Prompt) -- Public Data}
\begin{table}[H]
\centering
\caption{Evaluation Results for Llama 3.1 8B (With Prompt). $k = 0.9012$.}
\label{tab:8b_prompt}
\begin{tabular}{lcc}
\toprule
\textbf{Classification} & \textbf{Count} & \textbf{Rate (\%)} \\
\midrule
Abstain & 265 & 63.70\% \\
Faithful & 126 & 30.29\% \\
Inconsistent & 24 & 5.77\% \\
Hallucination & 1 & 0.24\% \\
\bottomrule
\end{tabular}
\end{table}
An interesting finding here is that the 8B model, despite being smaller, actually showed a slightly higher defensive rate (Abstain at 63.70\%) compared to the 17B model under the same System Prompt conditions. This suggests that the System Prompt works more effectively on smaller models: when encountering text with contradictions beyond their processing capacity, the 8B model tends to take the safe route of refusing to answer, rather than attempting the kind of complex reasoning that the 17B model goes for.

\subsection{Configuration 4: Llama 3.1 8B (No Prompt) -- Public Data}
\begin{table}[H]
\centering
\caption{Evaluation Results for Llama 3.1 8B (No Prompt). $k = 0.8056$.}
\label{tab:8b_noprompt}
\begin{tabular}{lcc}
\toprule
\textbf{Classification} & \textbf{Count} & \textbf{Rate (\%)} \\
\midrule
Abstain & 177 & 42.55\% \\
Inconsistent & 124 & 29.81\% \\
Faithful & 115 & 27.64\% \\
Hallucination & 0 & 0.00\% \\
\bottomrule
\end{tabular}
\end{table}
In sharp contrast to the collapse of the 17B No Prompt configuration, the 8B No Prompt model still maintained a reasonable Abstain rate (42.55\%). Interestingly, the model's more limited reasoning ability actually turned into an unexpected advantage here: when confronted with subtly mutated text full of contradictions, the model simply was not capable enough to reconcile the conflicting information, so it defaulted to Abstain out of confusion.

\subsection{Configuration 5: Llama 3.1 17B (With Prompt) -- Private Data}
\begin{table}[H]
\centering
\caption{Evaluation Results for Private Data. $k = 0.9180$.}
\label{tab:data_private}
\begin{tabular}{lcc}
\toprule
\textbf{Classification} & \textbf{Count} & \textbf{Rate (\%)} \\
\midrule
Abstain & 216 & 51.92\% \\
Faithful & 188 & 45.19\% \\
Inconsistent & 12 & 2.88\% \\
Hallucination & 0 & 0.00\% \\
\bottomrule
\end{tabular}
\end{table}
This configuration evaluates the 17B model on an internal (Private Data) dataset that is not publicly available online. Unlike the four previous configurations that used public data, the main goal here is to check whether the RAG system still hallucinates or ``corrects to the right answer'' (Inconsistent --- i.e., uses background knowledge to fix errors) when the model has no parametric knowledge about this type of data from pretraining.

The results show that the Inconsistent rate dropped sharply to just 2.88\% because the model simply does not have any relevant background knowledge to draw on for self-correction. On the other hand, the rate at which the model was fooled and passively accepted the misinformation (Faithful) jumped to 45.19\%. Notably, the Kappa score reached the highest value in the entire study (0.9180), demonstrating very high accuracy and consistency of the same AI Judge when applied to this dataset.

\subsection{Research Question Analysis}

\subsubsection{RQ1: Model Scale vs Vulnerability (Inverse Scaling Law)}

\textbf{Hypothesis:} We set out to test whether larger models become more vulnerable to semantic mutations when no defensive System Prompt is deployed.
\begin{itemize}
  \item $H_0$: The Abstain rate of the 17B model is greater than or equal to that of the 8B model (Abstain$_{17B}$ $\ge$ Abstain$_{8B}$). That is, the larger model is at least as safe as the smaller one.
  \item $H_1$: The Abstain rate of the 17B model is strictly lower than that of the 8B model (Abstain$_{17B}$ $<$ Abstain$_{8B}$). That is, the larger model is more vulnerable.
\end{itemize}

\textbf{Result:} The experimental data decisively \textbf{rejects $H_0$}. Under identical No Prompt conditions, the 17B model's Abstain rate collapsed to just 3.61\% (Table~\ref{tab:17b_noprompt}), while the 8B model maintained an Abstain rate of 42.55\% (Table~\ref{tab:8b_noprompt}). A Two-Proportion Z-test confirmed this difference is statistically significant ($p < 0.05$, $\alpha = 0.05$).

\textbf{Analysis:} This result confirms the \textit{Inverse Scaling Law of Semantic Poisoning}. The 17B model's superior instruction-following and reading comprehension abilities become a liability when unprotected: it seamlessly integrates flawed logic from the poisoned context and confidently produces Faithful responses (51.44\%), or detects the errors but overrides the context using its own parametric knowledge (Inconsistent: 44.71\%). The 8B model, lacking the cognitive capacity to reconcile contradictory information, defaults to the safer behavior of refusing to answer. In short, greater model capability amplifies vulnerability to semantic poisoning in the absence of explicit grounding instructions.

\subsubsection{RQ2: Data Domain \& Parametric Shield}

\textbf{Hypothesis:} We tested whether the model's ability to self-correct (Inconsistent Rate) diminishes when evaluated on Private Data that was not seen during pretraining.
\begin{itemize}
  \item $H_0$: The Inconsistent rate on Private Data is greater than or equal to that on Public Data (Inconsistent$_{Private}$ $\ge$ Inconsistent$_{Public}$). That is, the model's self-correction ability is maintained regardless of data domain.
  \item $H_1$: The Inconsistent rate on Private Data is strictly lower than that on Public Data (Inconsistent$_{Private}$ $<$ Inconsistent$_{Public}$). That is, the model loses its self-correction ability on unfamiliar data.
\end{itemize}

\textbf{Result:} The experimental data \textbf{rejects $H_0$}. On Private Data, the Inconsistent rate dropped to just 2.88\% (Table~\ref{tab:data_private}), compared to 9.37\% on Public Data in the equivalent 17B With Prompt configuration (Table~\ref{tab:17b_prompt}). A Two-Proportion Z-test confirmed this difference is statistically significant ($p < 0.05$, $\alpha = 0.05$).

\textbf{Analysis:} This confirms the \textit{Parametric Shield Hypothesis}. When processing Public Data (FastAPI documentation), the 17B model could cross-reference the mutated context against knowledge acquired during pretraining to detect inconsistencies. On Private Data, this ``parametric shield'' is entirely absent --- the model has no prior knowledge to fall back on. As a result, it becomes significantly more susceptible to blindly accepting poisoned information, with the Faithful rate surging to 45.19\%. This finding carries serious implications for enterprise RAG deployments: systems handling proprietary or internal data are inherently more vulnerable to semantic attacks than those operating on publicly available information.

